

The first line of input contains the number $N$, the length of the string $s$. This number will always be 
equal to $10^5$, except for in the sample. Your program does not have to solve the sample to get accepted.

The second line contains the string $s$ of length $N$. This string has exactly
$N/2$ characters \verb|b|, and their positions are generated uniformly at random.

Then, your program should start interacting with the grader. When it is grandpa's turn, you should read a 
single integer $i$ ($1 \leq i \leq N-3$), meaning that grandpa picks the berries at positions $i$, $i+1$, $i+2$, and $i+3$. When
it is your turn, you should instead print a number $i$, meaning that you pick the berries at positions $i$, $i+1$, $i+2$, and $i+3$.

If it is your turn and there are no more berries, your program should terminate without printing anything more.
If it is grandpa's turn and there are no berries, then the grader will not make any more moves, and your program should 
also terminate without printing anything more. If you picked at least as many berries as
grandpa, you pass the test case.

Each turn, grandpa greedily makes a move that picks as many berries as possible, and among all such moves, he chooses 
one uniformly at random.

The strings $s$ were generated beforehand, so they are the same across 
different submissions. However, grandpa's behaviour is randomized every submission, so it is 
possible to get different results if you submit the same code twice. 

There will be $100$ test cases. It is guaranteed that there exists a solution that passes this many cases with very high probability.

